\section{硅外延层的多光束识别与厚度反演}
\label{sec:silicon-results}

\subsection{反演范围与有效模型的具体实现}
附件3、4分别对应同一硅晶圆的 $10^\circ$、$15^\circ$ 入射光谱，各有7469个有效数值行，未发现重复波数。
两表首点均在 $399.6747\,\mathrm{cm}^{-1}$ 处取零，而下一点反射率跃升，故不以该端点参与条纹估计。
更主要的选窗依据是文献硅色散的适用范围：本文在
$450\leq\sigma\leq4000\,\mathrm{cm}^{-1}$ 内拟合，每角保留7363点，
联合样本量为14726。该窗口保留原始反射率，包括窄线形扰动；平滑曲线只用于辅助定位峰，不替代反演观测。

采用式\eqref{eq:silicon-drude-effective}--\eqref{eq:silicon-observation-model}。
外延层折射率固定为 Chandler-Horowitz 与 Amirtharaj 的室温参照色散\cite{ChandlerHorowitz2005}，
衬底由 $\Omega_p,\Gamma$ 描述其相对于参照的有效复介电响应。
对每个入射角分别计算 $s$、$p$ 偏振反射振幅，再作
$R_{\mathrm{Airy}}^{(\%)}=100R_{\mathrm{Airy}}=50(|r_s|^2+|r_p|^2)$ 的非偏振强度平均，数值单位为百分比。
随后以 $a_\theta+b_\theta R_{\mathrm{Airy}}^{(\%)}$ 拟合观测。
这一平均没有新增可调的偏振混合权重；非偏振近似是资料未提供偏振状态时的条件假设。

有效传播因子中的 $\exp[-g(\sigma/1000)^2]$ 对每次往返振幅施加相同的频率依赖衰减，
并同时进入 Airy 分子与分母，不能仅在最终反射率上乘一个包络。
代码中的 \texttt{eta} 对应本文 $g$。这一项使高波数条纹的振幅下降可被低维模型描述，
但不能将其唯一解释为吸收系数、粗糙度或相干长度；也没有把数据拟合成任意比例的相干、非相干反射率混合。
同样，$\Omega_p,\Gamma$ 只是本次反演的有效参数，不据此报告载流子浓度。

\subsection{变量投影与可复现的搜索设置}
共享非线性参数为 $\boldsymbol p=(d,\Omega_p,\Gamma,g)$。
每给定一组 $\boldsymbol p$，先计算两角理论谱，再通过各自的二元线性最小二乘
精确求得 $a_\theta,b_\theta$，最后最小化联合残差平方和。
这样以四维有界非线性搜索代替八维同时搜索，且保留两角各自的标定差异。
原始反射率百分点残差等权；最终两个 $b_\theta$ 均为正。

搜索边界设置为
\[
 1.5\leq d\leq6\;(\mu\mathrm m),\quad
 100\leq\Omega_p\leq15000,\quad
 1\leq\Gamma\leq5000\;(\mathrm{cm}^{-1}),\quad 0\leq g\leq2.
\]
这些是数值搜索范围，不是材料参数的置信区间。
主初值为 $(3.4,3900,370,0.10)$；另取
$d\in\{2.8,3.2,3.6,4.0,4.4\}$ 与
$\Omega_p\in\{3000,4500\}$ 的十组组合，
令 $\Gamma=350,g=0.08$，共进行11次有界最小二乘并选择残差最小解。
每次最多700次函数评价，三类停止容差均为 $10^{-10}$。
多起点确实发现了较差的局部极小点，因此只运行一个初值不足以保障结果。
最优解各参数均未触及边界，归一化雅可比矩阵条件数为4.22；
$d$ 与 $\Omega_p$ 的局部相关系数为0.854，表明固定色散后仍需关注相位与界面响应的耦合。

\subsection{计算结果与两束模型的公平比较}
最优共享参数为
\[
 \widehat d=3.405214\,\mu\mathrm m,\quad
 \widehat\Omega_p=3858.7745\,\mathrm{cm}^{-1},\quad
 \widehat\Gamma=367.7955\,\mathrm{cm}^{-1},\quad
 \widehat g=0.102003.
\]
标定参数分别为
$(a_{10},b_{10})=(-0.119354,0.937448)$、
$(a_{15},b_{15})=(-0.649952,1.047064)$；
$a_\theta$ 的单位为反射率百分点。
表\ref{tab:si-parameters}整理了参数及单位，标定系数在计算输出中按增益、偏置顺序保存。
\begin{table}[htbp]\centering
\caption{硅共同反演的参数记录}\label{tab:si-parameters}
\begin{tabular}{lrl}
\toprule
参数 & 估计值 & 单位\\\midrule
$d$ & 3.405214 & $\mu$m\\
$\Omega_p$ & 3858.7745 & $\mathrm{cm}^{-1}$\\
$\Gamma$ & 367.7955 & $\mathrm{cm}^{-1}$\\
$g$ & 0.102003 & 1\\
$a_{10}$ & -0.119354 & 百分点\\
$b_{10}$ & 0.937448 & 1\\
$a_{15}$ & -0.649952 & 百分点\\
$b_{15}$ & 1.047064 & 1\\
\bottomrule
\end{tabular}\end{table}

全文将上述多位小数用于复核，最终厚度建议表述为约 $3.41\,\mu\mathrm m$。

对照模型保留完全相同的材料响应、衰减、数据窗和标定参数，仅把无限反射级数截断为
\[
 r^{(2)}=r_{01}+(1-r_{01}^{\,2})r_{12}z_{\mathrm{eff}}.
\]
其中透射乘积 $1-r_{01}^{\,2}$ 被保留，避免把任意正弦回归误称为物理两束模型。
对照也采用同样的多起点设置并重新优化全部参数。
表\ref{tab:silicon-models}显示，多光束模型在相同参数数量下显著降低残差，
两束模型难以同时重建低波数尖谷与峰顶。
两种模型的共同厚度相差 $0.036336\,\mu\mathrm m$，约为多光束厚度的1.07\%；
这是当前数据和模型下的实际偏移，不意味着多光束在所有情形都会产生同方向的厚度偏差。

\begin{table}[htbp]
\centering
\small
\caption{硅光谱的模型对照；RMSE单位为反射率百分点}
\label{tab:silicon-models}
\begin{tabular}{lrrr}
\toprule
模型 & 参数数 & 共同厚度/$\mu\mathrm m$ & 联合RMSE\\
\midrule
两束，衬底有效响应与包络 & 8 & 3.368878 & 4.363185\\
多光束，不含额外包络（$g=0$） & 7 & 3.389471 & 1.201031\\
多光束，衬底有效响应与包络 & 8 & 3.405214 & 0.678945\\
\bottomrule
\end{tabular}
\end{table}

\paperfigure{figures/silicon_fit.png}{硅实测光谱与多光束共同厚度拟合。低波数尖谷为主要判别区域，高波数振幅逐渐下降。}{0.78}{fig:silicon-fit}

若强制 $g=0$，多光束模型的厚度为 $3.389471\,\mu\mathrm m$，
联合 RMSE 增至1.201031个百分点。这一消融说明有效包络能改善拟合，
也提醒厚度依赖于未充分测定的衰减模型，不能仅依靠最小残差宣布获得绝对真值。

\subsection{分角一致性、残差与外推能力}
表\ref{tab:silicon-angles}对照共享厚度与独立分角反演。
两个独立估计相差 $0.027718\,\mu\mathrm m$，约占共享厚度的0.81\%。
分角拟合明显减小残差，说明共享模型仍存在角间系统差异。
该差异可能涉及测点位置、波数标定、偏振或有效光学模型的不足；
现有附件不能将其归因于其中一种原因。
图\ref{fig:silicon-residuals}中的连续振荡残差也不支持把所有谱点视为独立同分布的白噪声。

\begin{table}[htbp]
\centering
\small
\caption{硅光谱的分角一致性检验}
\label{tab:silicon-angles}
\begin{tabular}{lrrr}
\toprule
估计方式 & 厚度/$\mu\mathrm m$ & $10^\circ$ RMSE & $15^\circ$ RMSE\\
\midrule
两角共享参数 & 3.405214 & 0.689061 & 0.668675\\
$10^\circ$独立反演 & 3.420549 & 0.229571 & 不参与拟合\\
$15^\circ$独立反演 & 3.392831 & 不参与拟合 & 0.326045\\
\bottomrule
\end{tabular}
\end{table}

\paperfigure{figures/silicon_residuals.png}{两束与多光束共同厚度模型的残差比较。残差仍有连续结构，条件拟合精度不能替代角间一致性。}{0.78}{fig:silicon-residuals}

为避免随机分点验证因相邻谱点高度相关而过于乐观，
将波数轴划分为宽度 $100\,\mathrm{cm}^{-1}$ 的连续块，末块按实际剩余范围取值，
同一块的两个角度同时划入同一折。
按块编号模5进行五折验证：每折用其余块重新估计全部参数，
再预测留出块，不使用留出观测重新标定。
多光束模型五折 RMSE 依次为
$0.8641,0.6457,0.4996,0.5838,0.8384$ 个百分点；
两束模型为
$6.5246,5.7939,4.7959,5.9071,3.3623$ 个百分点。
36个留出块中有34块支持多光束模型。
该结果表明改善并非仅源于样本内追随尖谷；它检验的是本数据波段内的分块预测，不能代替新晶圆或新仪器上的外部验证。
\paperfigure{figures/silicon_cv.png}{相同光学与标定结构下的五折连续块验证。每折均重新拟合模型，纵轴为留出数据的反射率 RMSE。}{0.76}{fig:silicon-cv}

\subsection{多光束证据的强度及其边界}
利用拟合模型计算有效往返保持因子
$\eta_{\mathrm{rt,eff}}=|r_{01}r_{12}z_{\mathrm{eff}}|$。
在 $10^\circ$ 下，$450$--$800$、
$800$--$1200$、$1200$--$2000$、
$2000$--$3000$、$3000$--$4000\,\mathrm{cm}^{-1}$
五个波段的中位数依次为
$0.3604,0.2668,0.0681,0.0162,0.0043$；
$15^\circ$ 的对应值非常接近，如表\ref{tab:si-roundtrip}所示。
\begin{table}[htbp]\centering
\caption{硅模型内的分波段平均往返保持量级}\label{tab:si-roundtrip}
\begin{tabular}{lrr}
\toprule
窗口/$\mathrm{cm}^{-1}$ & $10^\circ$中位数 & $15^\circ$中位数\\\midrule
450--800 & 0.36045 & 0.36040\\
800--1200 & 0.26675 & 0.26675\\
1200--2000 & 0.06811 & 0.06811\\
2000--3000 & 0.01620 & 0.01620\\
3000--4000 & 0.00425 & 0.00425\\
\bottomrule
\end{tabular}\end{table}

上述数值先对两种偏振的保持因子取平均。在前一波段，相对尾项量级指标
$\eta_{\mathrm{rt,eff}}/(1-\eta_{\mathrm{rt,eff}})$ 的中位数约为0.564，
该值不是两偏振强度误差的严格上界，但提示高阶往返不能当作微小修正；高波数端则明显减弱。
\paperfigure{figures/silicon_q.png}{由有效模型计算的往返振幅保持因子。该曲线依赖拟合模型，是支持性诊断而非独立的相干长度测量。}{0.78}{fig:silicon-q}

作为辅助检查，固定拟合相位并使用相同三次基线、二次谐波包络，
在 $1500$--$2800\,\mathrm{cm}^{-1}$ 比较一阶与二阶谐波回归。
两角RMSE分别由0.1821和0.1970降至0.0401和0.0673个百分点，
窗口中心的二次与基频振幅比分别为3.58\%和3.76\%。
该检查的相位来自主模型，不能当成独立检验，也不能把回归系数直接等同界面反射率。
综合尖谷形态、同参数物理对照、分块预测及往返保持量级，
附件3、4对多光束有效模型提供了较强支持。
但附件未给出光谱分辨率、光束发散、光斑重叠和厚度均匀性，
因此必要条件中的相干性与空间重叠不能逐项由数据独立核验。
本文据此采用多光束模型进行条件厚度反演，不作超出附件证据的绝对物理确证。

\subsection{厚度结果应附带的精度说明}
硅的窗口替换结果位于约$3.40248$--$3.40960\um$，而分角差异更大。
固定光学结构下120次配对循环块重抽样给出
$[3.403055,3.407313]\um$的条件分位区间；人为改变折射率$\pm1\%$，
厚度却在约$3.37075$--$3.44039\um$之间变化。
因此，建议报告条件厚度约$3.41\um$，同时列出分角结果，
不把求解器输出的小数位数解释成绝对测量精度。

后文的可靠性检验统一列出两材料的窗口、角度、折射率及重抽样检验，
并核对Airy闭式与显式反射级数、零折射率差极限和无噪声参数恢复。
这些检查分别约束模型稳定性与程序实现；提高实际厚度准确度仍需补充
独立折射率或掺杂信息、重复测量以及同测点的仪器标定。
